% Data-access contribution: Mason; final AI-assisted condensation.
\subsection{Data and SELECT queries}
All identities and observations are fictional. Three readable synthetic PDFs
match the seed; URLs use \nolinkurl{example.invalid}. The Python loader and
\texttt{seed.sql} yield identical dumps. Every table is populated: three ETFs,
six identifiers, thirteen categories, five Complete disclosures and seventeen
entries; each Complete total is one.

Q1 joins geography entries to funds and sources (8 rows). Q2 ranks Financials
allocations within one fictional scheme edition (2). Q3 groups by snapshot
with \texttt{COUNT} and \texttt{SUM} to show coverage/totals (5). Q4 separates
observation, publication and retrieval dates (5). Q6 traces entries to category
context and retained capture (17).

Q5 returns five \emph{diagnostic candidates}: holding weight multiplied by
candidate geography allocation. For example, $0.6\times0.985=0.591$ is not
validated exposure. Look-through also requires explicit holding/disclosure
dates, coverage/basis checks and compatible category meanings. The query
retains dates and contexts instead of silently aggregating them.

\subsection{Writes and reproducibility}
INSERT, UPDATE and DELETE bind user values through parameters and use the
transaction validator. Creation saves a snapshot and entries atomically; batch edits validate the
combined result. Invalid totals or missing batch entries roll back every change.
Deletion must preserve an entry and any Complete-total requirement; successful
Partial deletion is demonstrated in Figure~\ref{fig:ui}.

After installing dependencies, \texttt{quickstart.py} builds/seeds a missing
database; \texttt{--serve} starts the UI. Later runs validate/reuse data;
only \texttt{--reset} replaces it, after a successful temporary build.
The 74 tests comprise 33 schema, 11 integration and 30 UI-route tests.
Six-query output, seed equivalence, source-PDF consistency and reuse/reset
were also verified. Synthetic checks validate behavior, not provider accuracy.
